\chapter{Conclusion}

The goal of this thesis can be split up into two parts: 1) The investigation of Gilbert damping and other parameters in SC-FM-SC heterostructures near the critical temperature $T_\textrm{C}$ using FMR, and 2) the fabrication of a device capable of FMR measurements under a DC-bias current. 

Regarding the first part, the investigation of the parameters $\alpha$, $\Delta H_0$, $M_\textrm{eff}$ and $H_\textrm{shift}$ can be deemed partially succesful. The results for the sample with $d_\textrm{Nb} = \SI{45}{nm}$ are in good agreement with the literature \cite{Jeon2019} and previous measurements done by Mangold.
The results for the sample with $d_\textrm{Nb} = \SI{30}{nm}$ however show a different behavior than expected. While the parameters $\alpha$, $\Delta H_0$ and $H_\textrm{shift}$ show a similar trend as the 45 nm sample, the parameter $M_\textrm{eff}$ shows a completely different behavior, which could be attributed to sample aging and degradation. Near the critical temperature, almost all parameters show a spike-like feature, which has not been reported in literature yet. This feature could be attributed to domain formation or vortex motion in the superconductor near $T_\textrm{C}$, disturbing the spin-transfer mechanism.

The second part of the goal, the fabrication of a device capable of FMR measurements under a DC-bias current, can be considered successful. The microfabrication process was established and a first device was fabricated. However, due to time constraints, no measurements could be performed on this device. Future studies could use this device to investigate the influence of a DC-bias current on the parameters $\alpha$, $\Delta H_0$, $M_\textrm{eff}$ and $H_\textrm{shift}$ in SC-FM-SC heterostructures.